Unfortunately, we do not know how to compute the above oracles in $\mathsf{EXPH}$. The reason is that when we diagonalize $N_i$ on input length $n$, the number of heavy entries $P(r)$ we need to fix is $2^{(n'+4)K}\cdot2^n=2^{\operatorname{poly}(n')}$. Since $n'$ might be exponentially large compared to $n$, this upper bound might be doubly exponential. Nevertheless, we show that under the assumption that
\[\mathsf{SAT}\in\mathsf{DTIME}[n\cdot\operatorname{polylog}(n)]\cap\mathsf{NC},\tag{*}\]
the oracles $P,Q$ in Theorem11 can be computed in $\mathsf{PSPACE}$.
We sketch some intuition before presenting the full proof of Theorem12. The bottleneck of putting $\mathcal O$ into $\mathsf{EXPH}$ is that we do not have enough space to store all the heavy entries $P(r)$ fixed by a probabilistic machine $N_i$, as there might be doubly-exponentially many of them. Therefore, whenever we simulate a machine $N_i$ and it asks a query $P(r)$, we have to compute from scratch whether $P(r)$ was fixed by a previous machine. If we have an oracle $\mathsf{Heavy}$ that given $r$, decides whether $P(r)$ was already fixed, then we can simulate $N_i$ in exponential time with a constant number of alternations (using Theorem24). However, there are $\operatorname{poly}(n)$ machines $N_j$ simulated before $N_i$, and each time we invoke $\mathsf{Heavy}$, we need to simulate these machines to see if any of them has already fixed $P(r)$. It follows that the simulation of $N_i$ actually requires $\operatorname{poly}(n)$ alternations and $2^{\operatorname{poly}(n)}$ time. Still, under a strong enough assumption such as (*), we can simulate these $\operatorname{poly}(n)$ alternations in $2^{\operatorname{poly}(n)}$ time.
\begin{proof}[Proof of Theorem12]
Consider the construction in Theorem11, with the only difference that we do not calculate the precise probabilities; instead, we compute their approximations in $\mathsf{PH}$. Since $\mathsf{SAT}\in\mathsf{NC}$, we have $\mathsf{EXP}=\mathsf{PSPACE}$ by padding, thus it suffices to construct the oracles in $2^{\operatorname{poly}(n)}$ time.
The most involved part of this proof is to compute the following function $\mathsf{Heavy}(n,\langle P,Q\rangle,i,b,x_{i,b},r)$, which indicates whether $P(r)$ is a heavy query of the probabilistic machine $N_i$ on input $x_{i,b}$. Here, the input of $\mathsf{Heavy}$ consists of the description of oracles $P$ and $Q$, an integer $i\leq n$, an input $x_{i,b}\in\{0,1\}^n$, and a query $r\in\{0,1\}^{(n_{i'})^4}$ (where $n_{i'}$ is the smallest element in the sequence $\{n_i\}$ such that $n_{i'}\geq n$). We may assume $|r|\leq2^n$ as otherwise $N_i$ could never query $P(r)$. The oracle description $\langle P,Q\rangle$ will contain the following information:
\begin{itemize}
\item We record a table of all entries in $Q$ that are fixed; there are at most $\sum_{m\leq n}2m+2^{6n}\leq2^{O(n)}$ such entries. For each entry, we also record a timestamp indicating when this entry is fixed.
\item We record all entries of $P$ up to the $n_{i'-1}$-th slice (but we do not include any information of the $n_{i'}$-th slice of $P$); the description length of $P$ is also at most $2^{O(n)}$.
\end{itemize}
Note that the input length of $\mathsf{Heavy}$ is at most $2^{O(n)}$. Instead of requiring $\mathsf{Heavy}$ to decide exactly whether $P(r)$ is a heavy query, we only require a $2$-approximation: if $P(r)$ is queried by $N_i(x_{i,b})$ w.p. at least $2^{-(n'+4)K}$ then $\mathsf{Heavy}$ returns $1$, while if $P(r)$ is queried w.p. at most $2^{-(n'+4)K}/2$ then $\mathsf{Heavy}$ returns $0$.
Let $k$ be the number of machine-input pairs we have simulated before $N_i(x_{i,b})$ (inclusive). That is,
\[k:=\sum_{m<n}(2m)+2i+b+1\leq O(n^2).\]
We will show that:
\paragraph{Claim46.} If $\mathsf{SAT}\in\mathsf{DTIME}[n\cdot\operatorname{polylog}(n)]$, then we can compute $\mathsf{Heavy}$ in deterministic $O(2^{k^2})$ time.
\begin{proof}
Consider simulating (a random computational branch of) $N_i(x_{i,b})$ while answering the oracle queries of $N_i$ accordingly.\footnote{That is, whenever $N_i$ asks a query $P(r)$ or $Q(x,r)$ that is fixed, we answer accordingly. Whenever $N_i$ asks a query $P(r)$ that is not fixed, if $|r|=(n')^4$ then we return $b$, otherwise we return $0$. Whenever $N_i$ asks a query $Q(x,r)$ that is not fixed, we return $0$.} Note that we have not recorded the $n_{i'}$-th slice of $P$ in our description $\langle P,Q\rangle$. As a consequence, whenever $N_i(x_{i,b})$ asks a query $P(r_q)$ where $|r_q|=n_{i'}$, we need to recursively call $\mathsf{Heavy}$ to decide whether this query was already fixed and which value it was fixed to. We enumerate every $n_{i'-1}<\tilde n\leq n$, every machine $N_j$ ($j\leq\tilde n$), and every bit $b'$. Suppose that we feed the machine $N_j$ with the input $x'_{j,b'}\in\{0,1\}^{n_{i'}}$.\footnote{We assume that there is an easy and deterministic way of assigning the inputs $x_{j,b'}$ for each pair $(j,b')$.} Using the timestamps recorded in the table $Q$, we can recover the state of the oracle $Q$ before the simulation of $N_j(x'_{j,b'})$; we call this oracle $Q'$. (The portion of oracle $P$ up to the $n_{i'-1}$-th slice remains the same.) Then, we call $\mathsf{Heavy}(\tilde n,\langle P,Q'\rangle,j,b',x'_{j,b'},r_q)$ to see if this query is a heavy query fixed by $N'_j$. If it is, then we return $P(r_q):=b'$; otherwise we search through the next machine. If the query $P(r)$ was not fixed before and was asked in this simulation. then we return $1$; otherwise we return $0$.
The above argument implies that $\mathsf{Heavy}$ can be computed recursively in the following sense. There is a machine $V_{\mathsf{Heavy}}$ that gets $\mathsf{input}:=(n,\langle P,Q\rangle,i,b,x_{i,b},r)$ and some randomness $z\in\{0,1\}^{2^{O(n)}}$ (denoting a random computational branch of $N_i(x_{i,b})$), runs in $2^{O(n)}$ time with $\operatorname{poly}(n)$ invocations of the $\mathsf{Heavy}$ oracle, and outputs $0$ or $1$. Let $\varepsilon:=2^{-(n'+4)/K}$, we want that $\mathsf{Heavy}(\mathsf{input})=1$ if $\Pr_z[V_{\mathsf{Heavy}}(\mathsf{input};z)=1]\geq\varepsilon$, and $\mathsf{Heavy}(\mathsf{input})=0$ if $\Pr_z[V_{\mathsf{Heavy}}(\mathsf{input};z)=1]\leq\varepsilon/2$.
We can use induction to show that $\mathsf{Heavy}$ can be computed in (deterministic) $O(2^{k^2})$ time. In particular, by the induction hypothesis, each call of the $\mathsf{Heavy}$ oracle made by $V_{\mathsf{Heavy}}$ on $\mathsf{input}$ can be computed in $O(2^{(k-1)^2})$ time. Therefore, $V_{\mathsf{Heavy}}$ runs in at most $2^{(k-1)^2}\cdot\operatorname{poly}(k)$ time. By Theorem24, there is a $\Sigma_5\mathsf{TIME}[2^{(k-1)^2}\cdot\operatorname{poly}(k)]$ time machine $\widehat{\mathsf{Heavy}}$ such that $\widehat{\mathsf{Heavy}}(\mathsf{input})=1$ if $\Pr_z[V_{\mathsf{Heavy}}(\mathsf{input};z)=1]\geq\varepsilon$ and $\widehat{\mathsf{Heavy}}(\mathsf{input})=0$ if $\Pr_z[V_{\mathsf{Heavy}}(\mathsf{input};z)=1]\leq\varepsilon/2$. By our hypothesis that $\mathsf{SAT}\in\mathsf{DTIME}[n\cdot\operatorname{polylog}(n)]$, there is a deterministic machine $\mathsf{Heavy}$ that runs in $2^{(k-1)^2}\cdot\operatorname{poly}(k)<2^{k^2}$ time and decides the same language as $\widehat{\mathsf{Heavy}}$.
\end{proof}
Now we show that given an integer $N$, it is possible to print the $N$-th slice of $P$ and $Q$ in deterministic $2^{\operatorname{poly}(N)}$ time. We simulate the stages $n$ for $n=1,2,\ldots$, where during each stage we need to diagonalize against probabilistic machines $N_1,N_2,\ldots$. If $n=n_i$ for some $i\in\mathbb N$, then we also need to diagonalize against $M_i$.
\begin{itemize}
\item For each $N_i$, we can list the set of its heavy queries of the form $Q(x,r)$ in deterministic $2^{\operatorname{poly}(n)}$ time with a $\mathsf{PH}^{\mathsf{Heavy}}$ oracle. Note that $\mathsf{Heavy}$ can be decided in time quasi-polynomial in its input length; also recall we assumed that $\mathsf{SAT}\in\mathsf P$. Thus we can enumerate the heavy queries of the form $Q(x,r)$ in deterministic $2^{\operatorname{poly}(n)}$ time (without additional oracles). We fix all these queries. (Note that we do not fix the heavy queries of the form $P(r)$; instead, we use the oracle $\mathsf{Heavy}$ to decide whether a query $P(r)$ is fixed.) Then we use the $\mathsf{PH}$ oracle to estimate the probability that $N_i(x_{i,b})$ outputs $1$. If the probability is small then we pick some $r$ such that $Q(x_{i,b},r)$ is unfixed, and fix this entry to $1$; otherwise we do nothing.
\item Before we diagonalize against $M_i$, we spend $2^{\operatorname{poly}(n)}$ time to retrieve the list of fixed entries in the $n^4$-th slice of $P$ from the oracle $\mathsf{Heavy}$. Then we simulate $M_i$, and for every unfixed query, we fix it to be $0$. Let $b\in\{0,1\}$ be the output of $M_i$, then we set every unfixed entry in the $n^4$-th slice of $P$ to be $1-b$.
\end{itemize}
It is easy to see that the above process runs in $2^{\operatorname{poly}(N)}$ time.
\end{proof}
